\chapter{Data Integration Architecture for Knowledge-Intensive Domains}
\label{chap:architecture}

This chapter also consolidates the architectural prototypes we implemented across the various works in my \phd.
Indeed, starting from this works allows for a bottom up approach taking notes of the feedbacks received for the implementations in the two domains of civil judgements and criminal investigations.

\todo{use cases. why useful. users. challenges motivation}
faceted search, question answering, document exploration, annotation editing?, graph ui

\todo{richiamare principi da intro}


% \paragraph{Main challenges addressed in this work}
% In this work, we address several key challenges that arise when integrating \ac{ai}-based systems in knowledge-intensive and legally relevant domains:
% \begin{enumerate}[leftmargin=*,label=\textit{Ch.~\arabic*:},ref=\textit{Ch.~\arabic*}]
%     \item \label{ch:scalability} \emph{Scalability.} Legal and investigative datasets may comprise thousands of heterogeneous documents, requiring solutions that can efficiently process and integrate large volumes of text and metadata.
%     \item \label{ch:heterogeneity} \emph{Heterogeneity.} The documents vary widely in structure and type---ranging from judgments and contracts to chat logs and structured records---which demands flexible models capable of representing diverse sources within a unified framework.
%     \item \label{ch:traceability} \emph{Traceability and verifiability.} Users must be able to verify the outcome of \ac{ai}-driven operations by tracing each piece of extracted or inferred information back to its source document~\cite[Art.~15]{ai_act}. Furthermore, without such mechanisms, user trust and accountability in \ac{ai}-assisted decision-making may be compromised~\cite{HAN2024121052legalasst}.
% \end{enumerate}

% \paragraph{Objectives}
% Building on these challenges, the central objective of this \phd thesis is to design an architecture for data integration that addresses the \acl{ia} needs of professionals in the legal domain, supporting both formal legal documents---such as judgments and acts---and the heterogeneous data produced in investigative contexts.  
% The proposed architecture is guided by the following principles:
\begin{enumerate}
    \item \textbf{Generalizability} to heterogeneous documents and new use cases, as certain investigations may require ad-hoc workflows.
    \item data model
    \item uris? for external krs
    \item kr and eventually local kb
    \item RAE from claw
    \item different extractors. also from other domains
    \item standard format gatenlp. we borrow gatenlp model

    \item \textbf{Traceability and Verifiability}, ensuring that the origin of any extracted information can always be verified. For instance, a user searching for documents mentioning a person must be able to read the exact mention within its context, while a \acl{qa} system should provide references to the text passages supporting its answers.
    \item keep track of source documents
    \item annotation-based. what and where


    \item \textbf{Error correction capability}, allowing users to identify and correct extraction or linking errors, thus improving system quality with a \acf{hitl}~\cite{klie-etal-2020-zero} validation approach.
    \item hitl
    \item study in claw
    \item dave modification,\dots
    \item pipelines (easier to correct all the steps wrt end-to-end ? check in all the thesis)
    \item merge clusters

    \item \textbf{Scalability} to large and continuously growing document collections.
    \item microservices and parallelization
    \item study in claw
    \item mongo db ? are we so strict in the data db?
    
    \item IA UIs for the different use cases
    
    \item support for the investigation message use case: navigate in messages (e.g., get messages sent the same day,...)
    \item concept of subdocument? abstract annotation as subdocument?? annotation come subdocument?
    \item esempio 1-chat-per-doc e one-msg-per-doc. architettura deve supportare entrambi?! oppure in questi casi meglio usare neo4j?
    \item oppure neo4j è solo l'implementazione del data model? messaggi = docs con metadati. chat può essere considerata come un documento? ha senso annotarla per intero?
    \item idea: simulare questo use case e come farlo con data model e architettura
\end{enumerate}
%
% \nilchallenge

prototype? DAVE + IKBP
\begin{itemize}
    \item 
\end{itemize}

\section{Data Model}
GateNLP~\cite{cunningham-etal-2002-gate,gatenlp}




\section{Architecture}